\section{Squared multiplicities and hook deficits}\label{sec:multiplicities}
We write $\chi_\alpha$ for the character of $V_\alpha$ and
$\langle f,g\rangle_{S_m}=(m!)^{-1}\sum_{x\in S_m}f(x)\overline{g(x)}$ for the normalized character inner product. The cost $C(\alpha)$ in \eqref{eq:recursion} measures how many copies can occur when colors divide a row into four parts. We first justify this interpretation, then bound the cost by boxes outside a thin hook.

\begin{lemma}[Color-string trace]\label{lem:color-trace}
Let $\alpha\vdash m$, let $Q\ge0$ be a nonzero operator on $V_\alpha$, and define the probability
\[
 u(g)=\frac{D_\alpha}{m!\Tr Q^2}
             |\Tr(Q\rho_\alpha(g))|^2\quad(g\in S_m).
\]
Let $U$ be its average action on strings of length $m$ over four colors. Then
\[
 \Tr U\le\langle\chi_\alpha^2,4^{\#\mathrm{cycles}}\rangle_{S_m}
 =\sum_{\boldsymbol\beta}(c^{\alpha}_{\boldsymbol\beta})^2=C(\alpha).
\]
\end{lemma}
\begin{proof}
Schur orthogonality normalizes $u$ to have mass one. Write
$V_\alpha\otimes\overline{V_\alpha}=\bigoplus_\nu V_\nu\otimes\mathbb C^{a_\nu}$, and let $P_\nu$ be its isotypic projection. Character orthogonality gives
\[
 \sum_g u(g)\chi_\nu(g)
 =\frac{D_\alpha\Tr((Q\otimes\bar Q)P_\nu)}
               {D_\nu\Tr Q^2}.
\]
The two commuting positive operators $Q\otimes I$ and $I\otimes\bar Q$ satisfy
\[
 Q\otimes\bar Q\le\tfrac12(Q^2\otimes I+I\otimes\bar Q^2).
\]
The partial trace of $P_\nu$ over either factor commutes with the irreducible action on the other. Its trace is $a_\nu D_\nu$, so it equals
$(a_\nu D_\nu/D_\alpha)I$. The preceding character average is therefore at most $a_\nu$.

A permutation fixes exactly $4^{\#\mathrm{cycles}(g)}$ color strings. This is a permutation character, so its expansion in irreducible characters has nonnegative integer coefficients. Apply the preceding inequality to that expansion. It yields the first asserted bound, since $a_\nu$ is the multiplicity of $\nu$ in $\chi_\alpha^2$.

Finally the color-string representation is the direct sum, over all color-class sizes $m_1+\cdots+m_4=m$, of
$\Ind_{S_{m_1}\times\cdots\times S_{m_4}}^{S_m}\one$.
Frobenius reciprocity identifies its inner product with $\chi_\alpha^2$ as the sum of the dimensions of the commuting algebras of the corresponding restrictions of $V_\alpha$. Those dimensions are the sums of squared multiplicities displayed in the statement.
\end{proof}

For a partition $\alpha\vdash m$, put
\[
 k(\alpha)=m-\max(\alpha_1,\alpha'_1),\qquad
 t_h(\alpha)=|\{(i,j)\in\alpha:i>h,\ j>h\}|\quad(h\ge1).
\]
The latter is the number of boxes outside the union of the first $h$ rows and first $h$ columns.

\begin{lemma}[The cost of restriction]\label{lem:restriction-cost}
There is an absolute $C$ such that, for every $\alpha\vdash m$ and integer $h\ge1$,
\begin{equation}\label{eq:restriction-cost}
 \log C(\alpha)\le t_h(\alpha)\log4
       +C(1+h)(1+\sqrt{k(\alpha)})\log(2m).
\end{equation}
\end{lemma}
\begin{proof}
The central estimate for a residual shape $\delta\vdash t$ is
\begin{equation}\label{eq:squared-lr}
 (c^{\delta}_{\gamma_1,\ldots,\gamma_4})^2
 \le\binom{t}{|\gamma_1|,\ldots,|\gamma_4|}\le4^t.
\end{equation}
To prove it, let the multiplicity on the left before squaring be $c$. Restriction and induction dimensions give, respectively,
\[
 c\prod_iD_{\gamma_i}\le D_\delta,
 \qquad
 cD_\delta\le
 \binom{t}{|\gamma_1|,\ldots,|\gamma_4|}\prod_iD_{\gamma_i}.
\]
Multiplying and cancelling the positive dimensions proves \eqref{eq:squared-lr}. It is essential here to combine one restriction estimate and one induction estimate.

Both $C(\alpha)$ and $t_h(\alpha)$ are invariant under conjugation, so orient $\alpha$ with its first row of length $m-k$, where $k=k(\alpha)$. Remove its first $h$ rows and then the first $h$ columns of the remainder. The residual straight diagram $\delta$ has size $t=t_h(\alpha)$. Deleting a first row leaves an interlacing partition, so reversing that deletion adds a horizontal strip. Deleting a first column gives the vertical analogue.

More explicitly, let $\ell\le2h$ be the number of nonempty removed
lines, let their lengths be $a_1,\ldots,a_\ell$, and let $\xi_j$ be
the trivial character of $S_{a_j}$ for a removed row and the sign
character for a removed column. Then
\[
 W=\Ind_{S_t\times S_{a_1}\times\cdots\times S_{a_\ell}}^{S_m}
       (V_\delta\otimes\xi_1\otimes\cdots\otimes\xi_\ell)
\]
contains $V_\alpha$ by the Pieri rules. Thus each multiplicity in a restriction of
$V_\alpha$ is bounded by the corresponding multiplicity in $W$.
We will use this comparison only for tuples $\boldsymbol\beta$
that occur in the original restriction of $V_\alpha$. For those
tuples, the first-row inequality in the Littlewood--Richardson rule says
$\alpha_1\le\sum_i\beta_{i,1}$, and hence
\begin{equation}\label{eq:tail-budget}
 \sum_i(|\beta_i|-\beta_{i,1})\le k.
\end{equation}
Every intermediate diagram in a Pieri chain ending at such a $\beta_i$ is contained in it, and has at most $k$ boxes below its first row. The number of partitions of any size at most $m$ satisfying that condition is at most
\[
 Q_m(k)=(m+1)^{1+2\lceil\sqrt{k}\rceil}.
\]
One way to see this is to choose the first-row length, and encode the tail by its row lengths and column lengths along a Durfee square of side at most $\sqrt k$. Padding the two lists to length $\lceil\sqrt k\rceil$ gives the displayed upper bound. For $k=0$, only the first-row length is needed.

Fix such a final tuple $\boldsymbol\beta$, and put
$H_{\boldsymbol\beta}=\prod_{i=1}^4S_{|\beta_i|}$.
Mackey's formula for $\Res_{H_{\boldsymbol\beta}}W$ indexes the decomposition by
allocations of the residual factor and the $\ell$ strip factors
among these classes. There are at most $(m+1)^{4(\ell+1)}$
allocation matrices. For each allocation, restrict $V_\delta$ to
the four allocated subsets, then add the allocated horizontal or
vertical strips within each color class by Pieri. A chain ending
at $\beta_i$ consists of diagrams contained in $\beta_i$, so every
one obeys the small-tail bound just counted. Choosing the residual
tuple and all subsequent diagrams costs at most
$Q_m(k)^{4(\ell+1)}$. Each Pieri transition is multiplicity free,
and \eqref{eq:squared-lr} bounds the residual multiplicity by
$4^{t/2}$. Therefore
\[
 c^\alpha_{\boldsymbol\beta}
 \le \operatorname{mult}_{\boldsymbol\beta}(\Res_{H_{\boldsymbol\beta}}W)
 \le (m+1)^{4(\ell+1)}Q_m(k)^{4(\ell+1)}4^{t/2}.
\]
There are at most $Q_m(k)^4$ possible final occurring tuples.
Squaring this bound and summing over them gives $4^t$ times
$\exp\{C(1+h)(1+\sqrt k)\log(2m)\}$, proving
\eqref{eq:restriction-cost}. No tail estimate for the other
constituents of $W$ is required.
\end{proof}

The next fact lets us add hook costs along an entire level of the recursion.

\begin{lemma}[Hook deficits under restriction]\label{lem:hook-deficit}
If $c^{\lambda}_{\alpha_1,\ldots,\alpha_4}>0$, then for every integer $h\ge1$,
\begin{equation}\label{eq:hook-deficit}
 \sum_i t_h(\alpha_i)\le t_h(\lambda),
 \qquad \sum_i k(\alpha_i)\le k(\lambda).
\end{equation}
\end{lemma}
\begin{proof}
For the assertion about $k$, orient the parent with its longest side as first row. The first-row inequality gives a bound by its row deficit on the sum of the child row deficits; each child's $k$ is no larger than its row deficit. If the parent's longest side is a column, transpose every partition, which preserves Littlewood--Richardson multiplicities.

Write $s_\nu$ for the ordinary Schur function, and $s_\nu(X\mid Y)$ for its hook-Schur specialization \cite[Definition~1(a)]{orellana-zabrocki2000}. For the assertion about $t_h$, take $h$ ordinary variables $X$, $h$ transpose variables $Y$, and an unlimited ordinary alphabet $Z$. Variables in $X,Y$ have degree zero and those in $Z$ have degree one. The positive hook-Schur tableau rule uses ordinary entries weakly along rows and strictly down columns, and transpose entries strictly along rows and weakly down columns \cite[Section~2]{mason-niese2016}. For $s_\nu(X\mid Y)$, order all $X$ letters before all $Y$ letters, so the ordinary entries occupy a subdiagram. The rule gives
$s_\nu(X\mid Y)\ne0$ exactly when $\nu_{h+1}\le h$.
For completeness, a tableau over $h$ ordinary and $h$ transpose letters cannot fill an $(h+1)$-by-$(h+1)$ square: the ordinary subdiagram has at most $h$ rows, and the transpose-filled remainder has at most $h$ columns in each remaining row. Conversely, fill the first $h$ rows with their ordinary row letters and the remaining boxes with their transpose column letters. This realizes every diagram in that hook.

Writing $p_j$ for the $j$th power sum, the supersymmetric substitution is
\[
 p_j\longmapsto p_j(X)+(-1)^{j-1}p_j(Y).
\]
Applied to the ordinary Schur coproduct, it gives
\[
 s_\lambda(X+Z\mid Y)
 =\sum_{\nu\subseteq\lambda}s_\nu(X\mid Y)s_{\lambda/\nu}(Z).
\]
Its coefficients are nonnegative. The largest subdiagram of $\lambda$ in the $(h,h)$ hook is the union of its first $h$ rows and first $h$ columns. The smallest degree in $Z$ is therefore exactly $t_h(\lambda)$; the corresponding skew Schur function is nonzero on an unlimited alphabet.

Apply this substitution to
$\prod_i s_{\alpha_i}=\sum_\lambda c^{\lambda}_{\boldsymbol\alpha}s_\lambda$.
The product on the left has minimum $Z$-degree $\sum_i t_h(\alpha_i)$. Every summand on the right has nonnegative coefficients, so an occurring parent's minimum degree cannot be smaller. This is \eqref{eq:hook-deficit}.
\end{proof}

Equations~\eqref{eq:recursion} and \eqref{eq:restriction-cost} now give a recursive estimate in which the leading charge is $\theta\log4$ times a hook deficit. The other charges involve only $(1+h)(1+\sqrt k)\log(2m)$. We next prove where this recursion can stop, before summing those charges.
